\documentclass[11pt]{article}
\usepackage[a4paper,margin=26mm]{geometry}
\usepackage[T1]{fontenc}
\usepackage{lmodern,microtype}
\usepackage{amsmath,amssymb,amsthm,mathtools}
\usepackage{booktabs,array,enumitem}
\usepackage{xcolor}
\usepackage[colorlinks=true,linkcolor=blue!45!black,citecolor=blue!45!black,urlcolor=blue!45!black]{hyperref}
\hypersetup{pdftitle={Negative-Energy Coherent Dynamics: Radiation, Robust Profile Obstructions, and Mixed Dispersion},pdfauthor={Shuxing Yang},pdfsubject={Mass-critical dispersive dynamics: analytic research continuation, 18 October 2026}}
\usepackage{fancyhdr}
\pagestyle{fancy}\fancyhf{}
\fancyhead[L]{\small Negative-energy coherent dynamics}
\fancyhead[R]{\small Coherent dynamics}
\fancyfoot[C]{\thepage}
\setlength{\headheight}{14pt}
\setlength{\emergencystretch}{2em}
\setlist[enumerate]{itemsep=3pt,topsep=5pt}
\newtheorem{theorem}{Theorem}[section]
\newtheorem{proposition}[theorem]{Proposition}
\newtheorem{lemma}[theorem]{Lemma}
\newtheorem{corollary}[theorem]{Corollary}
\theoremstyle{definition}
\newtheorem{assumption}[theorem]{Profile input}
\newtheorem{definition}[theorem]{Definition}
\theoremstyle{remark}
\newtheorem{remark}[theorem]{Remark}
\DeclareMathOperator{\Rea}{Re}
\DeclareMathOperator{\Ima}{Im}
\DeclareMathOperator{\supp}{supp}
\DeclareMathOperator{\diag}{diag}
\DeclareMathOperator{\diver}{div}
\newcommand{\R}{\mathbb R}
\newcommand{\C}{\mathbb C}
\newcommand{\N}{\mathbb N}
\newcommand{\ip}[2]{\langle #1,#2\rangle}
\newcommand{\norm}[1]{\lVert #1\rVert}
\newcommand{\mc}{M_{\mathrm c}}
\newcommand{\etr}{M_{\mathrm{tr}}}
\newcommand{\ep}{\varepsilon}
\newcommand{\dd}{\,\mathrm d}
\title{\bfseries Negative-Energy Coherent Dynamics:\\
Radiation, Robust Profile Obstructions,\\and Mixed Dispersion}
\author{Shuxing Yang}
\date{18 October 2026}
\begin{document}
\maketitle
\begin{abstract}
We develop three extensions of a research program on negative-energy coherent structures in mass-critical focusing dispersive equations. First, using the existing four-dimensional excited-wave and spectral certificates as explicit inputs, we construct an exact forward-global solution consisting of one negative-energy biharmonic traveling wave and momentum-cancelling radiation. For a prescribed core with speed $v$ and momentum $-q e_1$, the infimum radiation mass compatible with zero total momentum and negative total energy is exactly $q(4/v)^{1/3}$; the construction attains every strictly larger radiation mass. Second, a finite-dimensional bordered reduction and scaling covariance yield a moment-free quantitative profile obstruction near the critical mass, uniformly in frequency and velocity, for general real-even homogeneous elliptic dispersion. Third, the phase-coercivity mechanism gives a mass gap for a genuinely three-dimensional cubic equation with one quadratic and two quartic dispersion directions, uniformly over every nonnegative quadratic correction in the latter directions. A joint extension makes this mixed-dispersion gap robust under profile residuals as well. Sharp Fourier-moment identities provide additional energy and parameter restrictions. The near-critical classification of general zero-momentum trajectories, the minimum traveling-wave mass endpoint, and existence for the new mixed-dispersion model remain open.
\end{abstract}

\section{A common problem and the new conclusions}

Negative energy is a central quantity in focusing evolution, but its consequences depend on dispersion and on the available coherent motion. In mass-critical biharmonic NLS, the existing project construction supplies excited traveling waves with negative energy and bounded Sobolev norms for all time. The companion mass-gap argument shows that these exact profiles cannot approach the sharp critical mass. These are complementary facts: a coherent mechanism exists, but critical variational geometry excludes it near the threshold.

The present continuation strengthens three parts of this picture. It reduces the coherent mass needed to obtain zero total momentum, makes the near-critical obstruction robust under profile error, and transfers the mechanism to a mixed-order model whose critical scaling is anisotropic. The connecting structure is the relation between energy, momentum, scaling, and the phase sector of the linearized action. Radiation tests how conserved quantities can be carried outside a core; a residual estimate tests how closely a general state can imitate a core; mixed dispersion tests which features survive a change of spatial scaling.

This manuscript contains new analytic derivations and their source dependencies. In particular, the first construction below is conditional on the specific existing project profile and spectral inputs stated in Section~\ref{sec:input}. The general profile obstruction uses the previously proved sharp-optimizer foundations, stated explicitly in Section~\ref{sec:robust}. The mixed-dispersion argument supplies those foundations directly in its concrete setting.

\subsection{Conventions}

The Fourier transform is unitary. Complex function spaces are treated as real spaces when taking variational derivatives. For a homogeneous symbol $a$ of degree $m>1$, set $s=m/2$, $p=2+2m/N$, and
\begin{equation}\label{eq:generalcharges}
 M(U)=\norm{U}_2^2,\qquad
 A(U)=\int a(\xi)|\widehat U(\xi)|^2\dd\xi,\qquad
 E(U)=\frac12A(U)-\frac1p\norm{U}_p^p.
\end{equation}
The momentum is $P(U)=\int\xi|\widehat U(\xi)|^2\dd\xi$. It is well defined in $H^s$ because $m>1$. A traveling wave $e^{i\omega t}U(x-vt)$ solves
\begin{equation}\label{eq:generalprofile}
 F(U;\omega,v):=a(D)U+\omega U+iv\cdot\nabla U-|U|^{p-2}U=0.
\end{equation}
Its action is $S_{\omega,v}=E+\omega M/2-v\cdot P/2$. For exact profiles the critical identities are
\begin{equation}\label{eq:generalpohozaev}
 v\cdot P=2mE,\qquad
 \omega M=\frac mN A+(2m-p)E.
\end{equation}
The sharp mass $\mc$ is defined by
\begin{equation}\label{eq:sharpGN}
 \norm{f}_p^p\le\frac p2\left(\frac{M(f)}{\mc}\right)^{m/N}A(f).
\end{equation}
Mass refers throughout to the squared $L^2$ norm, not to a relativistic particle mass.

\section{One coherent core and momentum-cancelling radiation}
\label{sec:radiation}

Here the equation is
\begin{equation}\label{eq:4dbnls}
 iu_t=\Delta^2u-|u|^2u,\qquad x\in\R^4,
\end{equation}
so $E(u)=\norm{\Delta u}_2^2/2-\norm{u}_4^4/4$.

\subsection{The existing profile and spectral inputs}\label{sec:input}

\begin{assumption}\label{ass:cert}
For every sufficiently small fixed $v>0$, the existing project certificates provide a smooth exponentially localized profile $Q_v$ solving
\[
 \Delta^2Q_v+Q_v+iv\partial_1Q_v-|Q_v|^2Q_v=0,
\]
with
\begin{equation}\label{eq:wavecharges}
 P(Q_v)=-q_v e_1,\qquad q_v>0,\qquad
 E(Q_v)=-\frac{vq_v}{8},\qquad
 q_v=2|I|v+O(v^3),\quad I<0.
\end{equation}
The branch tends to a real stationary excited state $Q$, and $M(Q_v)=M(Q)+O(v^2)$. It extends smoothly to frequency and vector-velocity parameters and obeys critical scaling. Its real action Hessian and Hamiltonian generator are
\[
 L_vh=(\Delta^2+1+iv\partial_1)h-2|Q_v|^2h-Q_v^2\bar h,
 \qquad B_v=-iL_v.
\]
The known spectral counts are
\begin{equation}\label{eq:spectralinput}
 n_-(L_v)=36,\qquad
 \dim E^u(B_v)=\dim E^s(B_v)=31,
\end{equation}
and $\ker L_v$ consists of the phase and four translation directions. The Hamiltonian invariant decomposition and its standard finite-dimensional hyperbolic estimates apply.
\end{assumption}

The original four-dimensional profile certificate gives $I<-175$ in Proposition~7.1 of the main construction~\cite{ProjectProfile}. The round-4 dynamics manuscript, Sections~2--3, supplies the Morse count 36, the five-dimensional symmetry kernel, and the positive charge Hessian with five length-two chains~\cite{ProjectR4}. The later continuation establishes the extra radial quartet and the exact unstable count in Sections~6.1--6.2~\cite{ProjectR5}. These are existing project inputs. The exact count 31 holds for sufficiently small nonzero speed; no explicit speed cutoff for that count is asserted here. The charge-Hessian identities used below also appear in the round-10 inertia section~\cite{ProjectR10}.

\begin{theorem}[A single core with zero total momentum]\label{thm:radiation}
Assume Project input~\ref{ass:cert}, and fix an admissible $v>0$. For every
\begin{equation}\label{eq:mucrit}
 \mu>\mu_{\rm crit}(v):=q_v(4/v)^{1/3},
\end{equation}
there are $r_+\in\mathcal S(\R^4)$, $T>0$, and an exact solution of \eqref{eq:4dbnls} in $C([T,\infty);H^{10})$ such that
\begin{equation}\label{eq:asymptotic}
 u(t)=e^{it}Q_v(\cdot-vte_1)+e^{-it\Delta^2}r_++o_{H^{10}}(1),
\end{equation}
and
\begin{equation}\label{eq:exactcharges}
 M(u)=M(Q_v)+\mu,\qquad P(u)=0,\qquad E(u)<0,
 \qquad \sup_{t\ge T}\norm{u(t)}_{H^{10}}<\infty.
\end{equation}
Time translation gives a forward-global solution from time zero, with the corresponding constant phase, spatial shift, and shifted scattering state in \eqref{eq:asymptotic}.

Among all solutions with the fixed core $Q_v$ and a strong $H^2$ asymptotic decomposition into that core and free radiation, zero total momentum and negative total energy require \eqref{eq:mucrit}. Hence the infimum total mass in this specified class is exactly
\begin{equation}\label{eq:classinf}
 M(Q_v)+q_v(4/v)^{1/3},
\end{equation}
and the infimum is not attained with negative total energy. In particular there are such solutions with $M(u)=M(Q)+O(v^{2/3})$ as $v\downarrow0$.
\end{theorem}

The previous symmetric two-wave construction has mass $2M(Q_v)$ and finite variance~\cite{ProjectR5}. Theorem~\ref{thm:radiation} uses one core and a radiation mass tending to zero. It does not approach the ground-state mass $\mc$: the limiting reference mass is that of the excited state $Q$. Finite variance and an open set of stable initial data are not conclusions of the new theorem.

\subsection{A center estimate from saturation of the negative index}

\begin{lemma}\label{lem:center}
Under Project input~\ref{ass:cert}, the center evolution satisfies
\begin{equation}\label{eq:center}
 \norm{e^{tB_v}P_c}_{H^{2+4k}\to H^{2+4k}}
 \le C_{v,k}(1+|t|),\qquad k\in\N\cup\{0\}.
\end{equation}
The stable and unstable evolutions satisfy exponential bounds in their respective decaying time directions.
\end{lemma}
\begin{proof}
Write $d(\omega,v)=S_{\omega,v}(U_{\omega,v})$. Critical scaling and the negative response give
\[
 d(\omega,v)=\omega d(1,v\omega^{-3/4}),\qquad
 d(1,v)=d(1,0)-\frac I2|v|^2+O(|v|^4).
\]
At $(\omega,v)=(1,ve_1)$ the leading axial Hessian is
\[
 \begin{pmatrix}-3Iv^2/8&Iv/2\\Iv/2&-I\end{pmatrix},
 \qquad \det=I^2v^2/8>0,
\]
and the transverse entries are $-I+O(v^2)>0$. Thus $d''$ is positive definite for sufficiently small nonzero $v$. The parameter derivatives generate five generalized symmetry directions with form matrix $-d''$. They contribute five negative directions, and invertibility of $d''$ makes the symmetry Jordan chains exactly length two.

In the Hamiltonian decomposition, the paired stable--unstable space has energy inertia $(31,31)$ and is energy-orthogonal to the complementary invariant blocks. Its 31 negative directions, together with the five symmetry-parameter directions, exhaust $n_-(L_v)=36$. The remaining invariant center space has positive energy. The free symbol $|\xi|^4+1-v\xi_1$ is positive for small $v$; the localized potential is relatively compact. Excluding the finite-dimensional kernel therefore makes the positive form coercive in $H^2$. Its conservation gives a uniformly bounded evolution on this complement, whereas the length-two symmetry chains grow at most linearly. This proves \eqref{eq:center} for $k=0$. This argument uses the invariant decomposition of Hamiltonian index theory~\cite{LinZeng}, with its operator hypotheses supplied by the stated project input; finite Morse index alone would not suffice.

Ellipticity and smooth bounded coefficients give the graph-norm equivalence
\[
 \norm h_{H^{2+4k}}\asymp\sum_{j=0}^k\norm{B_v^jh}_{H^2}.
\]
The invariant projections commute with the generator. Applying the $H^2$ estimates to each graph component proves the result for all fixed $k$.
\end{proof}

\subsection{Exact radiation cost and separation of velocities}

For any free radiation with mass $\mu$ and momentum $q_v e_1$, Jensen's inequality gives
\begin{equation}\label{eq:jensen}
 \frac12\norm{\Delta r_+}_2^2
 \ge\frac{q_v^4}{2\mu^3}.
\end{equation}
Consequently $-vq_v/8+\norm{\Delta r_+}_2^2/2<0$ requires precisely \eqref{eq:mucrit}. Equality in Jensen would force a point Fourier measure; it is unavailable to a nonzero $L^2$ function.

Conversely, fix $\mu>\mu_{\rm crit}$, let $k=q_v/\mu$ and $\eta=4k^3/v<1$, and choose a real radial $\phi\in C_c^\infty(B(0,1))$ with $\norm\phi_2=1$. For sufficiently small $\theta>0$ set $\sigma=\theta k$ and
\begin{equation}\label{eq:packet}
 \widehat r_+(\xi)=\mu^{1/2}\sigma^{-2}
 \phi\bigl((\xi-ke_1)/\sigma\bigr).
\end{equation}
Then $M(r_+)=\mu$, $P(r_+)=\mu k e_1=q_ve_1$, and
\[
 \frac12\norm{\Delta r_+}_2^2=\frac{\mu k^4}{2}C_\theta,
 \quad C_\theta=1+3\theta^2\int|\zeta|^2|\phi|^2
                  +\theta^4\int|\zeta|^4|\phi|^2.
\]
Choose $\theta$ so that $C_\theta\eta<1$ and $\beta=(1+\theta)^3\eta<1$. The total asymptotic energy is
\begin{equation}\label{eq:negbudget}
 -\frac{vq_v}{8}(1-C_\theta\eta)<0,
\end{equation}
and all radiative group velocities satisfy $|\nabla a(\xi)|\le\beta v<v$. The radiative packet and the core therefore separate by at least $(1-\beta)vt$ at the approximation level constructed next.

For a fixed choice $\eta\in(0,1)$ one has
\[
 \mu=q_v(4/\eta v)^{1/3}
 =2|I|(4/\eta)^{1/3}v^{2/3}+O(v^{8/3}).
\]
For a monochromatic packet of group speed $w=4k^3$, the energy cost is $q_vw/8$, and the combined energy is $q_v(w-v)/8$. The energy condition and separation condition thus agree at leading order.

\subsection{A localized high-order radiative approximation}

The free radiation alone leaves a cubic defect of order $t^{-4}$ in $H^{10}$. A center Green operator with linear growth would produce a borderline $t^{-2}$ correction. We remove that issue by constructing a radiative approximation to arbitrarily high order; no preservation of Fourier support by the exact nonlinear flow is assumed.

Put $y=x/t$, solve $y=4|\xi|^2\xi$ for $\xi=\xi(y)$, and let $\Phi(y)=y\cdot\xi(y)-|\xi(y)|^4$. The velocity image of \eqref{eq:packet} is compact and avoids zero. Define
\begin{equation}\label{eq:wkb}
 r_{\rm app}(t,x)=t^{-2}e^{it\Phi(x/t)}
                 \sum_{j=0}^Jt^{-j}A_j(x/t),\qquad
 A_0(y)=-\frac{\widehat r_+(\xi(y))}
                   {\sqrt{\det D^2a(\xi(y))}}.
\end{equation}
The minus sign is the four-dimensional stationary-phase factor. All coefficients are extended by zero outside a fixed compact velocity set.

For clarity, the recursion can be written explicitly. In the $y$ variable set
\[
 C_0=|\xi(y)|^2,\quad
 C_1=-i(2\xi(y)\cdot\nabla_y+\diver_y\xi),\quad C_2=-\Delta_y,
\]
and
\[
 \mathcal L_2=C_1^2+C_0C_2+C_2C_0,\quad
 \mathcal L_3=C_1C_2+C_2C_1,\quad \mathcal L_4=C_2^2.
\]
The order $t^{-1}$ term $C_0C_1+C_1C_0=-i(y\cdot\nabla_y+2)$ cancels the transport of the $t^{-2}$ prefactor. After eikonal cancellation, the amplitude equation is
\[
 i\partial_t A-t^{-2}\mathcal L_2A-t^{-3}\mathcal L_3A
 -t^{-4}\mathcal L_4A+t^{-4}|A|^2A=0.
\]
For $n\ge1$, set
\begin{equation}\label{eq:recursion}
 -inA_n-\mathcal L_2A_{n-1}-\mathcal L_3A_{n-2}
 -\mathcal L_4A_{n-3}
 +\sum_{a+b+c=n-3}A_a\overline{A_b}A_c=0,
\end{equation}
with negative-index coefficients and empty sums zero. Every step is a differential polynomial in previous coefficients divided by a nonzero integer. Hence all $A_j$ are smooth and retain the compact velocity support.

For every fixed Sobolev order $\ell$, differentiation and the change of variables $x=ty$ yield
\begin{align}
 \sup_{t\ge1}\norm{r_{\rm app}(t)}_{H^\ell}&<\infty,
 &\norm{r_{\rm app}(t)}_{W^{\ell,\infty}}&\le C_\ell t^{-2},\label{eq:wkbbounds}\\
 \norm{i\partial_tr_{\rm app}-\Delta^2r_{\rm app}
                +|r_{\rm app}|^2r_{\rm app}}_{H^\ell}
 &\le C_{\ell,J}t^{-J-2}.&&\label{eq:wkbdefect}
\end{align}
The physical factor $t^{-2}$ and the four-dimensional $L^2$ volume factor $t^2$ cancel. Spatial derivatives add bounded phase frequencies or negative powers of $t$. Stationary phase on the nondegenerate compact support, and integration by parts away from its velocity image, give
\begin{equation}\label{eq:freecompare}
 \norm{r_{\rm app}(t)-e^{-it\Delta^2}r_+}_{H^\ell}=O(t^{-1}).
\end{equation}
The support of $r_{\rm app}$ lies inside $|x|\le\beta vt$, while the core is centered at $vte_1$. Exponential localization of the core makes every mixed cubic defect exponentially small.

\subsection{The final-state argument and conserved quantities}

Set $R(t)=e^{it}Q_v(\cdot-vte_1)+r_{\rm app}(t)$. In the rotating and translating frame the correction $h=u-R$ obeys
\[
 \partial_t h=B_vh+f(t)+\mathcal N(t,h),
\]
where
\begin{align*}
 \norm{f(t)}_{H^{10}}&\le C(t^{-J-2}+e^{-ct}),\\
 \norm{\mathcal N(t,h)}_{H^{10}}
 &\le C\bigl[(t^{-4}+e^{-ct})\norm h_{H^{10}}
                       +\norm h_{H^{10}}^2+\norm h_{H^{10}}^3\bigr].
\end{align*}
The corresponding difference bound follows from the cubic algebra estimate. The $t^{-4}$ coefficient is the square of the radiative $W^{10,\infty}$ bound; terms containing both core and radiation are exponentially small.

Integrate center and unstable components backwards from infinity, and stable components forwards from a large time $T$ with zero stable datum. Lemma~\ref{lem:center} gives, for $L>2$,
\[
 \left\|\int_t^\infty e^{(t-\tau)B_v}P_cf(\tau)\dd\tau\right\|_{H^{10}}
 \le C_L\sup_{\tau\ge T}\tau^L\norm{f(\tau)}_{H^{10}}\,t^{2-L}.
\]
The hyperbolic integrals have exponentially localized kernels. On
\[
 X_J=\{h:\sup_{t\ge T}t^J\norm{h(t)}_{H^{10}}<\infty\},\qquad J\ge3,
\]
the linear radiation term has norm at most $CT^{-2}\norm h_{X_J}$ and the quadratic term at most $CT^{2-J}\norm h_{X_J}^2$. Choose the ball radius to contain the forced correction, then choose $T$ large. The map contracts and gives $h=O_{H^{10}}(t^{-J})$. Constants may depend on the fixed $v$ and $\mu$ and can diverge at a boundary of their admissible range.

The support separation, \eqref{eq:freecompare}, and $\norm{r_{\rm app}(t)}_4\to0$ yield asymptotic additivity of mass, momentum, and energy. Conservation therefore gives the exact identities
\[
 M(u)=M(Q_v)+\mu,\quad P(u)=P(Q_v)+P(r_+)=0,
 \quad E(u)=E(Q_v)+\tfrac12\norm{\Delta r_+}_2^2<0.
\]
Conversely, a strong $H^2$ decomposition of the stipulated class gives these same identities: free radiation tends to zero locally around any translated fixed profile and in $L^4$, by approximation with smooth dispersive states and the Sobolev embedding. Jensen's bound then proves necessity of \eqref{eq:mucrit}. This completes Theorem~\ref{thm:radiation}.

\section{A robust profile obstruction without spatial moments}\label{sec:robust}

Let $a$ be real, even, positive elliptic, $m$-homogeneous, and in the established smooth symbol class; for definiteness $a\in C^{N+2}(\R^N\setminus\{0\})$. Fix the charges \eqref{eq:generalcharges}. The following proof uses three earlier foundations: normalized sharp-GN near-optimizers are compact modulo translations and phase; every optimizer has constant phase; and for a real normalized optimizer $Q$, the phase form
\begin{equation}\label{eq:phaseform}
 L_-(Q)=a(D)+\omega_c-|Q|^{p-2},\qquad
 \omega_c=\frac{m}{N\mc},
\end{equation}
is nonnegative, has kernel $\R Q$, and is $H^s$-coercive on $Q^\perp$. These foundations are established in the preceding homogeneous-symbol work~\cite{ProjectR10}. They require neither real-sector nondegeneracy nor uniqueness of optimizers.

\begin{theorem}[Energy-controlled profile residual]\label{thm:robust}
There are $\delta>0$ and $C<\infty$ such that every $U\in H^s$ with
\[
 A(U)=1,\qquad \mc\le M(U)\le\mc+\delta,\qquad E(U)<0
\]
satisfies, for all $\omega\in\R$ and $v\in\R^N$,
\begin{equation}\label{eq:robustmain}
 |E(U)|\le C\norm{F(U;\omega,v)}_{H^{-s}}.
\end{equation}
The constants depend on the fixed symbol and dimension. No spatial moment is assumed, and the exponent of $|E|$ is optimal among uniform power estimates near zero energy.
\end{theorem}

\subsection{Parameter control and finite-dimensional reduction}

Sharp GN gives
\[
 -\frac12\left[(M(U)/\mc)^{m/N}-1\right]\le E(U)<0.
\]
Thus, as $\delta\to0$, normalized states in the theorem approach the compact normalized optimizer family, modulo translations and phase. Work in a chart $U=Q+h$ with $Q$ real. The map
\[
 (\alpha,b)\longmapsto\alpha Q+ib\cdot\nabla Q
 \quad\text{from }\R^{N+1}\text{ to }H^{-s}
\]
is injective: real and imaginary parts separate, and a nonzero $L^2$ function cannot be invariant in a nonzero translation direction. Since $m>1$, $\nabla:H^s\to H^{-s}$ is bounded. Subtracting $F(Q;\omega_c,0)=0$ and absorbing the parameter terms multiplying the small $h$ gives
\begin{equation}\label{eq:paramcontrol}
 |\omega-\omega_c|+|v|
 \le C_Q\bigl(\norm{F(U;\omega,v)}_{H^{-s}}+\norm h_{H^s}\bigr).
\end{equation}
A finite cover of the normalized optimizer set makes the constants uniform. If the residual exceeds a fixed small $\rho_0$, \eqref{eq:robustmain} follows at once from the GN energy bound. Hence we may work in a fixed small neighborhood of $(Q,\omega_c,0)$.

The map $F:H^s\to H^{-s}$ is real $C^1$. Indeed $p=2+2m/N$ lies strictly below the $H^s$ Sobolev endpoint, the map $L^p\ni u\mapsto |u|^{p-2}u\in L^{p'}$ is $C^1$ for every $p>2$, and dual Sobolev embedding completes the argument. Its derivative at the optimizer is
\[
 L=\diag(L_+,L_-),\qquad
 L_+=a(D)+\omega_c-(p-1)|Q|^{p-2}.
\]
Multiplication by $|Q|^{p-2}$ is compact from $H^s$ to $H^{-s}$: approximate it in $L^{p/(p-2)}$ by bounded compactly supported functions and use local Rellich compactness. Consequently $L$ is selfadjoint Fredholm of index zero and $K=\ker L$ is finite-dimensional.

Choose smooth compactly supported complex functions $\chi_1,\ldots,\chi_k$ whose real pairing matrix with a basis of $K$ is invertible. The bordered operator
\begin{equation}\label{eq:border}
 (h,c)\longmapsto
 \left(Lh+\sum_{j=1}^kc_j\chi_j,
                 (\ip h{\chi_j}_{\R})_{j=1}^k\right)
\end{equation}
is an isomorphism $H^s\times\R^k\to H^{-s}\times\R^k$. Projection onto $K$ determines $c$, the range equation determines $h$ modulo $K$, and the final constraints determine that kernel component.

The implicit function theorem therefore solves
\begin{equation}\label{eq:borderedsystem}
 F(V;\omega,v)+\sum_jc_j\chi_j=g,\qquad
 \ip{V-Q}{\chi_j}_{\R}=z_j
\end{equation}
as a $C^1$ function $(V,c)$ of $(\omega,v,z,g)$. For the original $U$, choose $z_j=\ip{U-Q}{\chi_j}_{\R}$. At $g=F(U)$ the solution is $(U,0)$. At $g=0$, denote it by $(V,c)$. With $\rho=\norm{F(U)}_{H^{-s}}$, local inverse bounds give
\begin{equation}\label{eq:regularize}
 \norm{U-V}_{H^s}+|c|\le C\rho,
 \qquad f:=F(V)=-\sum_jc_j\chi_j,\qquad \norm f_{H^{-s}}\le C\rho.
\end{equation}

\subsection{Scaling covariance supplies the missing regularity}

\begin{lemma}\label{lem:dilation}
The representative $V$ in \eqref{eq:regularize} satisfies
\[
 \Lambda V=(N/2+x\cdot\nabla)V\in H^s,
 \qquad \norm{\Lambda V}_{H^s}\le C
\]
uniformly in the fixed chart.
\end{lemma}
\begin{proof}
Let $T_\lambda f(x)=\lambda^{N/2}f(\lambda x)$. Homogeneity gives
\[
 F(T_\lambda V;\lambda^m\omega,\lambda^{m-1}v)
 =\lambda^mT_\lambda F(V).
\]
Thus $(T_\lambda V,c)$ solves \eqref{eq:borderedsystem} with parameters
\begin{align*}
 \omega_\lambda&=\lambda^m\omega,&v_\lambda&=\lambda^{m-1}v,\\
 z_{\lambda,j}&=\ip V{T_\lambda^{-1}\chi_j}_{\R}
                           -\ip Q{\chi_j}_{\R},&
 g_\lambda&=\sum_jc_j(\chi_j-\lambda^mT_\lambda\chi_j).
\end{align*}
All these parameters are $C^1$ in $\lambda$. Their derivatives involve the smooth test functions, not a derivative of $V$, and are bounded by a constant times $\norm V_2+|c|+|\omega|+|v|$. Initially $T_\lambda V$ is only strongly continuous in $H^s$, which suffices to place it in the local uniqueness neighborhood for $\lambda$ close to one. Uniqueness identifies it with the $C^1$ implicit-function solution along the displayed parameter curve. Hence it is $H^s$-differentiable. Uniform inverse bounds give the stated estimate. No moment of $V$ or $U$ was used in this argument.
\end{proof}

\subsection{The defect identity and completion of the estimate}

Rotate $V$ by a small constant phase so that $V=X+iY$ and $\ip YQ=0$. This preserves the residual norm and the dilation bound. Stable perturbation of \eqref{eq:phaseform} gives
\[
 \ip Y{B_VY}\ge c\norm Y_{H^s}^2,
 \qquad B_V=a(D)+\omega-|V|^{p-2}.
\]
The imaginary residual equation and the action derivative along the now legitimate $H^s$ scaling curve give, respectively,
\begin{equation}\label{eq:twoidentities}
 \ip Y{B_VY}=\frac12v\cdot P(V)+\ip Y{\Ima f},
 \qquad mE(V)-\frac12v\cdot P(V)=\Rea\ip f{\Lambda V}.
\end{equation}
Combining them, using Young's inequality and Lemma~\ref{lem:dilation}, yields
\[
 mE(V)=\ip Y{B_VY}-\ip Y{\Ima f}+\Rea\ip f{\Lambda V}
 \ge-C(\norm f_{H^{-s}}+\norm f_{H^{-s}}^2).
\]
Energy is locally Lipschitz on the bounded $H^s$ chart. Equation~\eqref{eq:regularize} transfers the bound to $U$ and proves Theorem~\ref{thm:robust}.

For arbitrary $A(U)>0$, scaling by $\lambda=A(U)^{-1/m}$ gives the equivalent formulation
\begin{equation}\label{eq:scaledresidual}
 |E(U)|\le C\left(\int
       \frac{|\widehat{F(U;\omega,v)}(\xi)|^2}
                    {1+a(\xi)/A(U)}\dd\xi\right)^{1/2}.
\end{equation}
Amplitude perturbations $(1+\ep)Q$, normalized back to $A=1$, have negative energy and a residual both of first order in $\ep$ at a suitable bounded frequency. Therefore an estimate $\rho\ge c|E|^\gamma$ with $\gamma<1$ is impossible uniformly near zero energy.

For an actual evolution $u(t,x)=e^{i\theta(t)}U(t,x-a(t))$, one has $F(U;\theta',a')=i\partial_tU$. The theorem is a lower bound on deformation speed modulo rigid phase and translation. It does not imply finite residence time: a curve may move indefinitely within a bounded set.

The previous residual--inertia theorem~\cite[Theorem 3.3, pp.~12--13]{ProjectR9} gives $v\cdot P\ge c|v|^2-C\rho^2$. Its stronger energy consequence $\rho^2\ge c_\theta(|E|+|v|^2)$ requires the additional dilation-defect assumption
\[
 D:=mE-v\cdot P/2\ge-\theta m|E|,\qquad 0\le\theta<1.
\]
Theorem~\ref{thm:robust} removes that assumption and proves the optimal linear bound for every state in the stated nearcritical class. An amplitude-increased real optimizer has $v=0$ and $D=-m|E|$, so it is outside every such restricted energy estimate and inside the present theorem. The two statements are complementary: the old bound has a stronger power on its restricted class, while the new argument enlarges the class without imposing a dilation defect or spatial moment. The present proof has received a separate internal mathematical check; external priority is not asserted.

\section{A uniform gap for three-dimensional mixed dispersion}\label{sec:mixed}

Consider one quadratic direction and two quartic directions,
\begin{equation}\label{eq:mixedpde}
 iu_t=-\partial_x^2u+\Delta_y^2u-\mu\Delta_yu-|u|^2u,
 \quad x\in\R,\ y\in\R^2,\ \mu\ge0.
\end{equation}
Write
\[
 A(U)=\norm{\partial_xU}_2^2+\norm{\Delta_yU}_2^2,
 \quad B_\mu(U)=\mu\norm{\nabla_yU}_2^2,
 \quad E_\mu(U)=\tfrac12(A+B_\mu)-\tfrac14\norm U_4^4.
\]
The principal energy space is $H_a=\{U\in L^2:A(U)<\infty\}$. The scaling $T_\lambda U(x,y)=\lambda^2U(\lambda^2x,\lambda y)$ preserves mass and scales $A$ by $\lambda^4$. Its homogeneous dimension is four, so the cubic nonlinearity is mass-critical for this anisotropic scaling. This is a three-variable model, not the isotropic fourth-order equation in three dimensions. Anisotropic critical exponents and blowup have an established literature~\cite{Fibich,Komada}; the assertion here concerns a uniform traveling-profile gap.

Define $\mc$ by the principal sharp inequality
\begin{equation}\label{eq:mixedGN}
 \norm U_4^4\le2\frac{M(U)}{\mc}A(U),
\end{equation}
and define the energy in the quadratic rest frame by
\begin{equation}\label{eq:internal}
 E_{\mu,\mathrm{int}}(U)=E_\mu(U)-\frac{P_x(U)^2}{2M(U)}.
\end{equation}

\begin{theorem}[Uniformity in the lower-order coefficient]\label{thm:mixed}
There exists a single $\delta>0$ such that, for every $\mu\ge0$, every frequency and velocity, and every nonzero weak $H_a$ profile solving
\begin{equation}\label{eq:mixedprofile}
 [-\partial_x^2+\Delta_y^2-\mu\Delta_y+\omega
       +iv_x\partial_x+iv_y\cdot\nabla_y]U-|U|^2U=0,
\end{equation}
one has
\[
 E_{\mu,\mathrm{int}}(U)<0\quad\Longrightarrow\quad
 M(U)\ge\mc+\delta.
\]
The same $\delta$ works for all nonnegative coefficients $\mu$, as well as for all frequencies and velocities. In particular the assertion holds if $E_\mu(U)<0$.
\end{theorem}

\subsection{Sharp compactness and phase coercivity}

The principal symbol $a(\xi,\eta)=\xi^2+|\eta|^4$ has sublevel volume proportional to $R^4$ on $\{a\le R^4\}$. Hausdorff--Young and weighted Cauchy--Schwarz give $H_a\hookrightarrow L^q$ for every finite $q\ge2$. Scaling this inhomogeneous bound gives \eqref{eq:mixedGN} with a finite positive best constant. Also $H_a\hookrightarrow H^1(\R^3)$, so local Rellich compactness in $L^2$, followed by interpolation, gives local compactness in $L^4$. The ordinary $H^1$ local-mass vanishing lemma, combined with an $L^q$ bound for $q>4$, rules out global $L^4$ mass when all translated local $L^2$ masses vanish.

Normalize a maximizing sequence by $M(U_j)\to\mc$, $A(U_j)=1$, and $\norm{U_j}_4^4\to2$. After translation there is a nonzero weak limit $Q$. Let $M_1=M(Q)$, $A_1=A(Q)$, $M_2=\mc-M_1$, and $A_2=1-A_1$. Quadratic splitting, Brezis--Lieb, and \eqref{eq:mixedGN} imply
\[
 \mc\le M_1A_1+M_2A_2.
\]
If both masses were positive, the right side would be strictly less than $\mc(A_1+A_2)=\mc$. Hence $M_2=0$, and the same inequality forces $A_2=0$. This proves attainment and normalized near-optimizer compactness.

The real and complex sharp constants agree. Indeed, for $U=X+iY$, the triangle inequality in the density space $L^2$ and Cauchy--Schwarz give
\[
 \norm{X^2+Y^2}_2\le\norm{X^2}_2+\norm{Y^2}_2
 \le(2/\mc)^{1/2}\bigl(M(X)^{1/2}A(X)^{1/2}
                              +M(Y)^{1/2}A(Y)^{1/2}\bigr).
\]
Equality at an optimizer forces $X^2$ and $Y^2$ to be proportional. Apply the same fact to $(X+Y,X-Y)/\sqrt2$; the second proportionality removes a spatially varying relative sign. Thus every optimizer has constant phase. Its real representative need not have one sign.

A normalized real optimizer obeys
\[
 (-\partial_x^2+\Delta_y^2+\omega_*)Q-Q^3=0,
 \qquad \omega_*=1/\mc.
\]
Let $L_Q=-\partial_x^2+\Delta_y^2+\omega_*-Q^2$. Expansion of the sharp inequality at $Q+ith$ proves $L_Q\ge0$. If $h\in\ker L_Q$, put $a=A(h)$ and $b=M(h)/\mc$. H\"older and sharp GN give
\[
 a+b=\int Q^2h^2\le\norm Q_4^2\norm h_4^2
 \le2\sqrt{ab}\le a+b.
\]
Equality implies $h^2=bQ^2$. Apply the same conclusion to $h+Q$, which also lies in the form kernel; subtraction gives $h=cQ$ with a constant real $c$. Multiplication by $Q^2$ is compact relative to $H_a$, by approximation in $L^2$ and local compactness. Consequently
\begin{equation}\label{eq:mixedcoercive}
 \ip h{L_Qh}\ge c_Q(M(h)+A(h)),\qquad \ip hQ=0.
\end{equation}
This supplies all optimizer ingredients used in the contradiction proof.

\subsection{Profile identities without a moment assumption}

A weak profile in $H_a$ lies in $H^1(\R^3)\cap L^6$, so its cubic nonlinearity and first-order drifts belong to $L^2$. The equation gives $(-\partial_x^2+\Delta_y^2)U\in L^2$, and the symbol controls the ordinary $H^2$ norm. Since $H^2(\R^3)$ is an algebra, elliptic bootstrapping with the full mixed operator gives $U\in H^k$ for every finite $k$.

Let $z=(x,y)$, $F_0(z)=(2x,y)$, and choose a real compact cutoff $\chi$ equal to one near zero. Define
\[
 F_R(z)=\chi(z/R)F_0(z),\quad d_R=\diver F_R,
 \quad K_R=F_R\cdot\nabla+d_R/2.
\]
This generator is skew-adjoint in $L^2$. Its differentiated coefficients satisfy $\norm{\partial^\alpha F_R}_\infty\le C_\alpha R^{1-|\alpha|}$ for $|\alpha|\ge1$. Inside the cutoff ball, $d_R=4$ and $DF_R=\diag(2,1,1)$.

For the quadratic term,
\[
 \partial_x(K_RU)=K_R\partial_xU+
 \sum_j(\partial_xF_{R,j})\partial_jU+\tfrac12(\partial_xd_R)U.
\]
The skew term has zero real pairing with $\partial_xU$, and the limit is $2\norm{\partial_xU}_2^2$. Similarly,
\begin{align*}
 \Delta_y(K_RU)={}&K_R\Delta_yU
 +2\sum_{\alpha\in y,j}(\partial_\alpha F_{R,j})
                         \partial_\alpha\partial_jU
 +(\Delta_yF_R)\cdot\nabla U\\
 &+(\nabla_yd_R)\cdot\nabla_yU+\tfrac12(\Delta_yd_R)U.
\end{align*}
Pairing with $\Delta_yU$ gives limit $2\norm{\Delta_yU}_2^2$. Annular errors are bounded by the $H^2$ tail, plus $R^{-1}$ and $R^{-2}$ times bounded lower norms, and tend to zero. The lower-order term gives $B_\mu$. The cubic term contributes $\int d_R|U|^4/4\to\norm U_4^4$, and the frequency term vanishes. The drift commutator gives limit $-v_xP_x-v_y\cdot P_y/2$. Therefore
\begin{equation}\label{eq:mixedpohoz}
 4E_\mu=B_\mu+v_xP_x+\tfrac12v_y\cdot P_y.
\end{equation}
Testing the equation with $U$ also gives
\begin{equation}\label{eq:mixednehari}
 A+B_\mu+\omega M-v_xP_x-v_y\cdot P_y-\norm U_4^4=0.
\end{equation}
Finally, a compactly generated $x$-phase variation with $\phi_R(x)=x\chi_1(x/R)$ yields
\[
 \int\phi_R'(x)\bigl(2\Ima(\bar U\partial_xU)-v_x|U|^2\bigr)=0.
\]
Dominated convergence gives $2P_x=v_xM$. All limiting arguments use Sobolev tails, not $xU$ or $yU$ in $L^2$.

\subsection{Quadratic rest frame and the uniform contradiction}

Put $U=e^{iv_xx/2}W$. Then $W$ solves the profile equation with $v_x=0$ and $\omega'=\omega-v_x^2/4$, and
\[
 P_x(W)=0,\qquad M(W)=M(U),\qquad
 E_\mu(W)=E_\mu(U)-\frac{P_x(U)^2}{2M(U)}.
\]
The reduced identities are
\begin{equation}\label{eq:reducedmixed}
 4E_\mu(W)=B_\mu+\tfrac12v_y\cdot P_y,
 \qquad \omega'M=A-B_\mu+4E_\mu.
\end{equation}
Suppose, contrary to the theorem, that reduced negative-energy profiles have masses decreasing to $\mc$, allowing arbitrary coefficients $\mu_j\ge0$. Normalize their principal kinetic energies to one. Under this scaling the coefficient becomes $\mu_{j,\lambda}=\lambda^2\mu_j$, so the class remains unchanged. Sharp GN gives
\[
 E_{\mu_j}(W_j)\ge\tfrac12B_j
                      -\tfrac12(M(W_j)/\mc-1).
\]
Thus $B_j\to0$, $E_{\mu_j}(W_j)\to0$, and $\norm{W_j}_4^4\to2$. Near-optimizer compactness gives $W_j\to Q$ strongly in $H_a$, after translation and phase, and \eqref{eq:reducedmixed} gives $\omega'_j\to1/\mc$.

Choose the remaining phase so that $Y_j=\Ima W_j\perp Q$. Strong $L^4$ convergence makes
\[
 -\partial_x^2+\Delta_y^2+\omega'_j-|W_j|^2
\]
uniformly coercive on $Q^\perp$. Adding $-\mu_j\Delta_y$ preserves this property because it is a nonnegative form. The imaginary profile equation and \eqref{eq:reducedmixed} give
\[
 4E_{\mu_j}(W_j)=B_j+
 \ip{Y_j}{(-\partial_x^2+\Delta_y^2-\mu_j\Delta_y+
                   \omega'_j-|W_j|^2)Y_j}\ge0,
\]
a contradiction. It is unnecessary to prove that the normalized $\mu_j$ themselves tend to zero. This proves Theorem~\ref{thm:mixed}.

The proof also permits any positive-semidefinite quadratic form in the $y$ variables. If $v_y=0$, \eqref{eq:reducedmixed} gives $4E_{\mu,\mathrm{int}}=B_\mu\ge0$; negative internal energy therefore requires motion in the quartic directions. The theorem does not establish the existence of such profiles in the mixed model.

\subsection{A joint robust obstruction in the quadratic rest frame}

The bordered argument also applies to the mixed model itself. This joins the preceding two mechanisms rather than treating them as separate assertions. Let $K$ be a symmetric positive-semidefinite $2\times2$ matrix and replace $B_\mu$ by
\[
 B_K(U)=\int\eta^{\mathsf T}K\eta\,|\widehat U(\xi,\eta)|^2\dd\xi\dd\eta,
 \qquad E_K(U)=\tfrac12(A+B_K)-\tfrac14\norm U_4^4.
\]
The scalar case is $K=\mu I$. Write $E_{K,\mathrm{int}}=E_K-P_x^2/(2M)$ and
\[
 R_{K,\omega,v}(V)=
 [a(D)+D_y^{\mathsf T}KD_y+\omega+iv_x\partial_x
              +iv_y\cdot\nabla_y]V-|V|^2V.
\]

\begin{corollary}[Uniform robust mixed-dispersion obstruction]\label{cor:robustmixed}
There are $\delta>0$ and $C<\infty$, independent of $K\ge0$, such that the following holds. For a nonzero state $U$, let
\[
 V=e^{-iP_x(U)x/M(U)}U.
\]
If $A(V)=1$, $\mc\le M(U)\le\mc+\delta$, and $E_{K,\mathrm{int}}(U)<0$, then
\begin{equation}\label{eq:robustmixed}
 |E_{K,\mathrm{int}}(U)|\le
 C\inf_{\omega,v_x,v_y}\norm{R_{K,\omega,v}(V)}_{H_a^*}.
\end{equation}
Here the dual norm uses Fourier denominator $1+\xi^2+|\eta|^4$. No spatial moment is assumed. The boost and the principal kinetic normalization are part of the statement.
\end{corollary}
\begin{proof}
The boost gives $P_x(V)=0$ and $E_K(V)=E_{K,\mathrm{int}}(U)$. Put $d=M(V)/\mc-1$. Principal GN implies
\[
 0\le B_K(V)<d,\qquad |E_K(V)|\le d/2,
 \qquad 2<\norm V_4^4\le2(1+d).
\]
Thus $V$ approaches the normalized principal optimizer family. For a real optimizer define the $y$-gradient covariance
\[
 (C_Q)_{ij}=\Rea\int\partial_{y_i}\bar Q\,\partial_{y_j}Q.
\]
It is positive definite: a null vector would make a nonzero $L^2$ function invariant along a spatial direction. Continuity in $H_a$ and compactness of the optimizer family give a uniform $C_Q\ge cI$. Consequently $C_V\ge cI/2$ near that family, and
\begin{equation}\label{eq:Ksmall}
 B_K(V)=\operatorname{tr}(KC_V)\ge(c/2)\operatorname{tr}K,
 \qquad \norm K\le Cd.
\end{equation}
Thus the originally arbitrary lower-order coefficient becomes uniformly small in the normalized nearcritical chart.

The parameter absorption of \eqref{eq:paramcontrol} now gives
\[
 |\omega-\omega_c|+|v|\le C(\rho+\norm{V-Q}_{H_a}+\norm K),
 \qquad \rho=\norm{R_{K,\omega,v}(V)}_{H_a^*}.
\]
The large-residual case follows from $|E_K(V)|\le d/2$. In the small-residual case, apply the bordered implicit function theorem with $K$ as an additional finite-dimensional parameter. It is legitimate to allow signed symmetric $K$ in the local inverse, while retaining $K\ge0$ in the assertion. The resulting representative $W$ has
\[
 \norm{W-V}_{H_a}\le C\rho,\quad f=R_{K,\omega,v}(W),
 \quad\norm f_{H_a^*}\le C\rho,\quad |P_x(W)|\le C\rho.
\]
The last estimate uses $P_x(V)=0$ and local Lipschitz continuity of momentum; the regularized $W$ need not itself be in the rest frame.

Under anisotropic dilation, the parameters scale as
\[
 K_\lambda=\lambda^2K,\quad\omega_\lambda=\lambda^4\omega,
 \quad(v_x)_\lambda=\lambda^2v_x,\quad(v_y)_\lambda=\lambda^3v_y.
\]
The same compact-test-function argument as Lemma~\ref{lem:dilation} proves
\[
 \Lambda_aW=(2+2x\partial_x+y\cdot\nabla_y)W\in H_a,
 \qquad\norm{\Lambda_aW}_{H_a}\le C.
\]
After a constant phase rotation write $W=X+iY$, $Y\perp Q$, and set
\[
 L_W=a(D)+D_y^{\mathsf T}KD_y+\omega-|W|^2.
\]
Its phase form is uniformly coercive on $Q^\perp$; the $K$ term is nonnegative. Combining the imaginary residual equation with the anisotropic action derivative at fixed $K$ gives the exact identity
\begin{align*}
 4E_K(W)={}&B_K(W)+\ip Y{L_WY}+\tfrac12v_xP_x(W)\\
 &-\ip Y{\Ima f}+\Rea\ip f{\Lambda_aW}.
\end{align*}
The first term is nonnegative, the second absorbs the imaginary-residual pairing, and the remaining terms are bounded below by $-C(\rho+\rho^2)$. Uniform local Lipschitz continuity of $E_K$, using \eqref{eq:Ksmall}, transfers this lower bound to $V$. The constant is uniform over the finite optimizer-chart cover and all $K\ge0$. Since the estimate holds for each frequency and velocity, taking their infimum proves \eqref{eq:robustmixed}.
\end{proof}

Setting the residual to zero recovers the mixed traveling-profile mass gap. In this form the exclusion theorem and the residual theorem are consequences of one quantitative mechanism in the quadratic rest frame.

\section{What has advanced, and the next decisive questions}

The three results address different parts of the same coherent-dynamics problem. The single-core construction realizes an exact zero-momentum negative-energy trajectory and computes the least radiation mass within a prescribed asymptotic class. The robust obstruction controls all nearcritical negative-energy states by their profile defect, uniformly over the moving frame. Corollary~\ref{cor:robustmixed} carries that same quantitative mechanism into anisotropic critical scaling with an arbitrary nonnegative lower-order coefficient, and includes the mixed exact-profile gap as its zero-residual consequence.

Three questions remain central. First, can a negative-energy, zero-momentum solution with mass arbitrarily close to $\mc$ remain globally bounded? The present radiation construction stays near the excited-state mass, and the robust residual estimate excludes rigid approximation rather than recurrent or genuinely time-dependent bounded motion. A theorem about all trajectories needs an additional dynamical mechanism.

Second, define $\etr$ as the infimum mass of negative-energy exact traveling profiles. The existing compactification gives either a strong negative-energy minimizer or a finite collection of zero-energy profiles at common parameters~\cite{ProjectR10}. The moment restrictions in Appendix~\ref{app:moments} improve normalization, but do not decide which alternative occurs. On a smooth parameter branch, a negative-energy mass minimum must be a charge-Hessian degeneracy. Neither the known negative directions nor the residual theorem excludes such a turning point.

Third, the mixed model has a uniform exclusion theorem but no negative-internal-energy construction in this continuation. Establishing a real excited profile with a strictly negative response in a quartic direction, followed by a justified branch continuation, would supply a concrete realization. This is a new existence problem, not a consequence of the exclusion argument. These questions keep the program centered on actual coherent structures, their mass thresholds, and their dynamical alternatives.

\appendix
\section{Sharp Fourier-moment restrictions and the mass endpoint}\label{app:moments}

The following identities supply unconditional parameter restrictions for pure BNLS profiles. Their constants are sharp for probability measures and for smooth Fourier densities. Sharpness among exact nonlinear profiles is not asserted.

\begin{lemma}[A sharp fourth-moment identity]\label{lem:momentone}
Let $\nu$ be a probability measure on $\R^N$ with finite fourth moment. Set $\mu=\int\xi\dd\nu$, $q=\int|\xi|^2\xi\dd\nu$, and $a_4=\int|\xi|^4\dd\nu$. Then
\begin{align}\label{eq:momentidentity}
 a_4+8\mu\cdot q
 =\int\Bigl[&(|\xi|^2+4\mu\cdot\xi-8|\mu|^2)^2\notag\\
 &+16\bigl(|\mu|^2|\xi|^2-(\mu\cdot\xi)^2\bigr)\Bigr]\dd\nu\ge0.
\end{align}
For an absolutely continuous probability measure the inequality is strict.
\end{lemma}
\begin{proof}
The integrand expands to $|\xi|^4+8(\mu\cdot\xi)|\xi|^2-64|\mu|^2(\mu\cdot\xi)+64|\mu|^4$. The last two terms cancel after integration. If $\mu\ne0$, equality forces support on the two points $(-2\pm2\sqrt3)\mu$; if $\mu=0$, equality forces support at zero. Neither is possible for an absolutely continuous probability measure.

For sharpness take $\mu=e$, $|e|=1$, and masses
\[
 p_+=\frac{3+2\sqrt3}{4\sqrt3},\qquad
 p_-=\frac{2\sqrt3-3}{4\sqrt3}
\]
at $(-2+2\sqrt3)e$ and $(-2-2\sqrt3)e$, respectively. Then $q=-8e$ and $a_4=64$. Smooth narrow Fourier packets approximate these two atoms, proving optimality of the coefficient $1/8$ in $\mu\cdot q\ge-a_4/8$.
\end{proof}

For a pure-BNLS profile in dimension $N$, use $\dd\nu=|\widehat U|^2\dd\xi/M$. The no-moment group-velocity identity is $v=4q$, while $v\cdot P=8E$. Lemma~\ref{lem:momentone} gives
\begin{equation}\label{eq:energybound}
 E(U)>-A(U)/16.
\end{equation}
In fact the integral in \eqref{eq:momentidentity}, multiplied by $M$, equals $A+16E$. In dimensions $2\le N<12$, a negative-energy profile therefore satisfies
\begin{equation}\label{eq:generalparameters}
 \frac{N\omega}{4}<\frac AM<\frac{8N\omega}{3(12-N)},
 \qquad -\frac{N\omega M}{6(12-N)}<E<0.
\end{equation}
This follows from $\omega M=(4/N)A+(6-8/N)E$. In particular
\[
 \begin{array}{c|c|c}
 N&\text{kinetic-to-mass range}&\text{energy range}\\\hline
 3&3\omega/4<A/M<8\omega/9&-\omega M/18<E<0\\
 4&\omega<A/M<4\omega/3&-\omega M/12<E<0.
 \end{array}
\]
At $N=12$, $\omega>0$ still follows strictly but the displayed upper bound on $A/M$ is unavailable. At $N=1$, negative energy instead gives $4A<\omega M<33A/8$.

\begin{lemma}[A second sharp moment bound]\label{lem:momenttwo}
With the same notation,
\begin{equation}\label{eq:momenttwo}
 a_4+2\mu\cdot q\ge\frac3{2^{2/3}}|q|^{4/3}.
\end{equation}
It is strict for absolutely continuous probability measures.
\end{lemma}
\begin{proof}
For $q\ne0$, let $e=q/|q|$ and $d=(2|q|)^{1/3}$. Direct expansion gives
\begin{align*}
 &\int\left[(|\xi|^2-d e\cdot\xi-d^2/2)^2
                  +d^2(|\xi|^2-(e\cdot\xi)^2)\right]\dd\nu\\
 &\hspace{3em}=a_4+2\mu\cdot q-3d^4/4\ge0.
\end{align*}
Equality forces support on $d(1\pm\sqrt3)e/2$. With respective weights $(\sqrt3-1)/(2\sqrt3)$ and $(\sqrt3+1)/(2\sqrt3)$, these atoms have $\mu=0$, $q=d^3e/2$, and $a_4=3d^4/4$. This also proves sharpness, by smooth approximation. The $q=0$ case is immediate.
\end{proof}

In four dimensions the exact identities give $\omega=a_4+2\mu\cdot q$. Thus every four-dimensional cubic BNLS profile, without an energy-sign assumption, satisfies
\begin{equation}\label{eq:cone}
 \omega>\frac3{2^{10/3}}|v|^{4/3},\qquad
 |v|<\frac{4\sqrt2}{3^{3/4}}\omega^{3/4}.
\end{equation}
This necessary condition is weaker than positivity of the full free moving symbol, which requires $\omega>3(|v|/4)^{4/3}$; no Fredholm gap is inferred from \eqref{eq:cone} alone.

In dimension three, a negative-energy profile has $\mu\cdot q<0$ and $\omega=\tfrac43(a_4+2\mu\cdot q)-\mu\cdot q$. Lemma~\ref{lem:momenttwo} therefore also gives the concise necessary bound $|v|<2\omega^{3/4}$.

A further dimension-independent bound holds for negative-energy profiles. With $a_2=\int|\xi|^2\dd\nu$, Cauchy--Schwarz gives $|q-a_2\mu|^2\le a_2(a_4-a_2^2)$. Since $\mu\cdot q<0$, maximization in $a_2$ implies
\[
 |q|^2<\frac2{3\sqrt3}a_4^{3/2},\qquad
 |v|<4(4/27)^{1/4}(A/M)^{3/4}.
\]

These restrictions allow the minimum traveling-wave mass problem in dimensions three and four to be posed at fixed $\omega=1$, with uniformly controlled kinetic-to-mass ratio and velocity. They do not settle attainment. More explicitly, if a negative-energy minimizer lies in a smooth full parameter family, then local mass minimality gives $M_\omega=D_vM=0$. For $d=S(U_{\omega,v})$, one has $d_\omega=M/2$, $d_v=-P/2$, hence $P_\omega=0$. Momentum scaling yields
\[
 (m-1)D_vP\,v=P,\qquad d''(\omega,v)(1,0,\ldots,0)^{\mathsf T}=0.
\]
For an isotropic quartic branch $P=p(c)e$, $c=|v|$, this requires $3cp'(c)=p(c)<0$. Controlling these actual charge-turning points, including possible singular branches, is the remaining endpoint issue.

\begin{thebibliography}{99}
\raggedright
\bibitem{ProjectProfile}
Project research manuscript, \emph{Nonradial negative-energy traveling waves for the mass-critical biharmonic NLS}, 7 October 2026. Source: \path{BNLS_Nonradial_Research_20261007.tex}; especially Proposition 7.1 and the traveling-branch construction in Section 8. Existing profile and interval-certificate input.

\bibitem{ProjectR4}
Project round-4 research continuation, \path{BNLS_Dynamics_20261007.pdf}, 7 October 2026, Sections 2--3. Existing Morse count, symmetry kernel, charge-Hessian sign, and generalized-kernel inputs; not recertified here.

\bibitem{ProjectR5}
Project research continuation, \emph{Zero-momentum global dynamics and three-dimensional traveling waves for the mass-critical biharmonic NLS}, 7 October 2026. Source: \path{BNLS_Zero_Momentum_and_Three_Dimensions_20261007.pdf}, Sections 1--3 and 6.1--6.2. Existing quartet transfer and exact unstable-count inputs; not recertified here.

\bibitem{ProjectR9}
Project round-9 research continuation, \path{BNLS_Dilation_Structure_and_Profile_Gaps_20261008.pdf}, 8 October 2026, Theorem 3.3, pp.~12--13. Quantitative residual--inertia inequality and its conditional energy consequence.

\bibitem{ProjectR10}
Project research continuation, \emph{BNLS thresholds, spectral geometry and radiation}, 8 October 2026. Proof sources: \path{round10_foundations_section.tex}, \path{round10_threshold_section.tex}, \path{round10_inertia_section.tex}, and \path{round10_independent_audit.md}. Includes the homogeneous-symbol phase foundations, threshold compactification, and charge Hessian.

\bibitem{LinZeng}
Z. Lin and C. Zeng, \emph{Instability, index theorem, and exponential trichotomy for linear Hamiltonian PDEs}, Memoirs of the American Mathematical Society \textbf{275} (2022), no. 1347. \url{https://arxiv.org/abs/1703.04016}.

\bibitem{CoteLeCoz}
R. C\^ote and S. Le Coz, \emph{High speed excited multi-solitons in nonlinear Schr\"odinger equations}, Journal de Math\'ematiques Pures et Appliqu\'ees \textbf{96} (2011), 135--166. \url{https://arxiv.org/abs/1007.3034}. Background on excited multi-solitons; not the present fourth-order final-state result.

\bibitem{Fibich}
G. Fibich, B. Ilan and S. Schochet, \emph{Critical exponents and collapse of nonlinear Schr\"odinger equations with anisotropic fourth-order dispersion}, Nonlinearity \textbf{16} (2003), 1809--1821. \href{https://doi.org/10.1088/0951-7715/16/5/314}{doi:10.1088/0951-7715/16/5/314}. Author manuscript: \url{https://www.math.tau.ac.il/~fibich/Manuscripts/anisotropic_NL.pdf}.

\bibitem{Komada}
K. Komada, \emph{Existence of blow-up solutions to nonlinear Schr\"odinger equations with anisotropic fourth-order dispersion}, Funkcialaj Ekvacioj \textbf{65} (2022), 191--214. \href{https://doi.org/10.1619/fesi.65.191}{doi:10.1619/fesi.65.191}.
\end{thebibliography}
\end{document}
